% SGA 3, Expose I, Loop 1 English reconstruction.
% Scope: Section 3 in full, from its heading through the paragraph immediately
% before Section 4; printed/component pages 18--20.
% Authority: Polo--Gille Exp1-13oct24.pdf, component pages 18--20
% (combined-reader pages 28--30). Force-OCR is drafting/locator material only.

\section*{\SGAInlineRefTarget{sga3:i:section:3}3. \(\mathbf O\)-模范畴与
\(\mathbf G\)-\(\mathbf O\)-模范畴}
\addcontentsline{toc}{section}{3. O-模范畴与 G-O-模范畴}

\medskip
\noindent\SGAInlineRefTarget{sga3:i:definition:3-1}\textbf{定义 3.1.}
设 \(\mathbf O\) 和 \(\mathbf F\) 是
\(\widehat{\mathscr C}\) 的两个对象. 若对每个
\(S\in\operatorname{Ob}(\mathscr C)\), 集合 \(\mathbf O(S)\) 配备了
一个环结构, 而 \(\mathbf F(S)\) 配备了这个环上的模结构, 并且对于
\(\mathscr C\) 中的每个箭头 \(S'\to S''\), 映射
\(\mathbf O(S'')\to\mathbf O(S')\) 都是环同态, 且
\(\mathbf F(S'')\to\mathbf F(S')\) 都是与该环同态相容的阿贝尔群同态,
则称 \(\mathbf F\) 是 \(\widehat{\mathscr C}\)-环 \(\mathbf O\)
上的一个 \(\widehat{\mathscr C}\)-模, 或简称为 \(\mathbf O\)-模.

固定 \(\mathbf O\) 后, 按通常方式定义两个 \(\mathbf O\)-模
\(\mathbf F\) 与 \(\mathbf F'\) 之间的态射, 从而得到阿贝尔群
\(\operatorname{Hom}_{\mathbf O}(\mathbf F,\mathbf F')\),
以及 \(\mathbf O\)-模范畴, 记作 \((\mathbf O\text{-Mod.})\).

\medskip
\noindent\SGAInlineRefTarget{sga3:i:lemma:3-1-1}\textbf{引理 3.1.1.}\footnotemark
\textit{范畴 \((\mathbf O\text{-Mod.})\) 具有逐对象定义的阿贝尔范畴结构.
此外, \((\mathbf O\text{-Mod.})\) 满足公理 \(\mathrm{(AB\ 5)}\)
(见 \([\mathrm{Gr}57,\,1.5]\)); 即任意直和都存在, 并且若
\(\mathbf M\) 是一个 \(\mathbf O\)-模, \(\mathbf N\) 是一个子模,
而 \((\mathbf F_i)_{i\in I}\) 是 \(\mathbf M\) 的一个递增滤过子模族, 则
\[
  \bigcup_{i\in I}(\mathbf F_i\cap\mathbf N)
    =\left(\bigcup_{i\in I}\mathbf F_i\right)\cap\mathbf N.
\]}

事实上, 设 \(f:\mathbf F\to\mathbf F'\) 是一个 \(\mathbf O\)-模态射.
对每个 \(S\in\operatorname{Ob}(\mathscr C)\), 置
\[
  (\operatorname{Ker}f)(S)=\operatorname{Ker}f(S)
\]
(分别地, 置相应的类似公式), 以此定义 \(\mathbf O\)-模
\(\operatorname{Ker}f\) (分别地,
\(\operatorname{Im}f\) 和 \(\operatorname{Coker}f\)).
则 \(\operatorname{Ker}f\) (分别地, \(\operatorname{Coker}f\))
是 \(f\) 的一个核 (分别地, 余核), 并且有 \(\mathbf O\)-模同构
\[
  \mathbf F/\operatorname{Ker}f
    \xrightarrow{\ \sim\ }\operatorname{Im}f.
\]
这就证明了 \((\mathbf O\text{-Mod.})\) 是阿贝尔范畴.

任意直和都存在, 并且是逐对象定义的. 最后, 若
\(\mathbf M\) 是一个 \(\mathbf O\)-模, \(\mathbf N\) 是一个子模,
而 \((\mathbf F_i)_{i\in I}\) 是 \(\mathbf M\) 的一个递增滤过子模族, 则包含关系
\[
  \bigcup_{i\in I}(\mathbf F_i\cap\mathbf N)
    \subset
  \left(\bigcup_{i\in I}\mathbf F_i\right)\cap\mathbf N
\]
实际上是等式. 事实上, 若 \(S\in\operatorname{Ob}(\mathscr C)\) 且
\[
  x\in\mathbf N(S)\cap\bigcup_i\mathbf F_i(S),
\]
则存在一个 \(i\in I\), 使得
\(x\in\mathbf N(S)\cap\mathbf F_i(S)\).

\medskip
\noindent\SGAInlineRefTarget{sga3:i:proposition:3-1-2}\textbf{命题 3.1.2.}\footnotemark[\value{footnote}]
\textit{假设范畴 \(\mathscr C\) 是小范畴, 即
\(\operatorname{Ob}(\mathscr C)\) 是一个集合.
则 \((\mathbf O\text{-Mod.})\) 有一个生成元, 因而有足够多的内射对象.}
\footnotetext{编者注: 新增\SGARefLink{sga3:i:lemma:3-1-1}{断言 3.1.1} 与
\SGARefLink{sga3:i:proposition:3-1-2}{3.1.2}, 是为了说明范畴
\((\mathbf O\text{-Mod.})\) 是阿贝尔范畴且满足公理
\(\mathrm{(AB\ 5)}\), 并且当范畴 \(\mathscr C\) 是小范畴时,
还有足够多的内射对象. 此处修正了 Linyuan Liu 于 2016 年指出的
\SGARefLink{sga3:i:proposition:3-1-2}{3.1.2} 中的一个错误.}

\textit{证明.}
对 \(\mathscr C\) 的每个对象 \(S\), 如下定义 \(\mathbf O\)-模
\(\mathbf O_S\). 对每个对象 \(S'\), 令 \(\mathbf O_S(S')\) 为
以 \(\operatorname{Hom}(S',S)\) 为指标的 \(\mathbf O(S')\) 的若干副本的直和, 即
\[
  \mathbf O_S(S')
    =\bigoplus_{g:S'\to S}\mathbf O(S')g.
\]
对每个态射 \(f:S''\to S'\), 以 \(f^*\) 表示环同态
\(\mathbf O(S')\to\mathbf O(S'')\), 则态射
\[
  \mathbf O_S(f):\mathbf O_S(S')\longrightarrow\mathbf O_S(S'')
\]
将 \(\mathbf O_S(S')\) 的每个元素 \(ag\), 其中
\(a\in\mathbf O(S')\) 且 \(g:S'\to S\), 映到 \(\mathbf O_S(S'')\)
的元素 \(f^*(a)(g\circ f)\). 容易验证, 这定义了一个取值于
\(\mathbf O\)-模的函子.

由于假设 \(\operatorname{Ob}(\mathscr C)\) 是一个集合, 可以构造直和
\[
  \mathbf G=\bigoplus_{S\in\operatorname{Ob}(\mathscr C)}\mathbf O_S.
\]

\medskip
\noindent\SGAInlineRefTarget{sga3:i:lemma:3-1-2-1}\textbf{引理 3.1.2.1.}
\textit{设 \(\mathbf F\) 是一个 \(\mathbf O\)-模.
对每个 \(S\in\operatorname{Ob}(\mathscr C)\), 每个元素
\(x\in\mathbf F(S)\) 都定义了唯一的 \(\mathbf O\)-模态射
\[
  \mathbf x:\mathbf O_S\longrightarrow\mathbf F
\]
使得 \(\mathbf x(1\operatorname{id}_S)=x\).}

\textit{证明.}
设 \(\phi:\mathbf O_S\to\mathbf F\) 是一个 \(\mathbf O\)-模态射,
满足 \(\phi(1\operatorname{id}_S)=x\).
设 \(g:S'\to S\) 且 \(a\in\mathbf O(S')\).
由 \(\mathbf O(S')\)-线性以及下图的交换性

% Authority: Exp1-13oct24.pdf component/printed p. 19; combined reader p. 29.
\begin{center}
  \begin{tikzcd}[column sep=4.5em,row sep=3.0em]
  \mathbf O_S(S)
    \arrow[r,"\phi"]
    \arrow[d,"\mathbf O_S(g)"']
  &
  \mathbf F(S)
    \arrow[d,"\mathbf F(g)"]
  \\
  \mathbf O_S(S')
    \arrow[r,"\phi"]
  &
  \mathbf F(S')
\end{tikzcd}

\end{center}

可得
\[
  \phi(ag)
    =a\phi(1g)
    =a\mathbf F(g)\bigl(\phi(1\operatorname{id}_S)\bigr)
    =a\mathbf F(g)(x).
\]
因此, 若 \(\phi\) 存在, 则它由上述等式确定.
反之, 若如此定义 \(\phi\), 则对每个 \(f:S''\to S'\), 都有
\begin{align*}
  \phi\circ\mathbf O_S(f)(ag)
    &=\phi\bigl(f^*(a)(g\circ f)\bigr)\\
    &=f^*(a)\mathbf F(g\circ f)(x)\\
    &=\mathbf F(f)\bigl(a\mathbf F(g)(x)\bigr)\\
    &=\mathbf F(f)\circ\phi(ag),
\end{align*}
即下图

% Authority: Exp1-13oct24.pdf component/printed p. 19; combined reader p. 29.
\begin{center}
  \begin{tikzcd}[column sep=4.5em,row sep=3.0em]
  \mathbf O_S(S')
    \arrow[r,"\phi"]
    \arrow[d,"\mathbf O_S(f)"']
  &
  \mathbf F(S')
    \arrow[d,"\mathbf F(f)"]
  \\
  \mathbf O_S(S'')
    \arrow[r,"\phi"]
  &
  \mathbf F(S'')
\end{tikzcd}

\end{center}

是交换的. 引理得证.

现在设 \(\mathbf F\) 是一个 \(\mathbf O\)-模.
对每个 \(S\in\operatorname{Ob}(\mathscr C)\), 令 \(F_0(S)\) 为
\(\mathbf O(S)\)-模 \(\mathbf F(S)\) 的一个生成元集, 并令
\(\mathbf L_S\) 为以 \(F_0(S)\) 为指标的 \(\mathbf O_S\) 的若干副本的直和.
由引理, 得到一个 \(\mathbf O\)-模态射
\[
  \mathbf L_S\longrightarrow\mathbf F
\]
使得态射 \(\mathbf L_S(S)\to\mathbf F(S)\) 是满射,
因而是 \(\mathbf O(S)\)-模范畴中的满态射.

置
\[
  I=\bigsqcup_{S\in\operatorname{Ob}(\mathscr C)}F_0(S).
\]
于是得到一个 \(\mathbf O\)-模态射
\[
  \mathbf G^{\oplus I}
    \longrightarrow
  \bigoplus_{S\in\operatorname{Ob}(\mathscr C)}\mathbf L_S
    \longrightarrow
  \mathbf F
\]
它逐对象为满态射, 因而是 \((\mathbf O\text{-Mod.})\) 中的满态射.
这表明 \(\mathbf G\) 是 \((\mathbf O\text{-Mod.})\) 的一个生成元
(见 \([\mathrm{Gr}57,\,1.9.1]\)). 由于
\((\mathbf O\text{-Mod.})\) 满足 \(\mathrm{(AB\ 5)}\),
由 \([\mathrm{Gr}57,\,1.10.1]\) 可知,
\((\mathbf O\text{-Mod.})\) 有足够多的内射对象.

\medskip
\noindent\SGAInlineRefTarget{sga3:i:remark:3-1-3}\textbf{注 3.1.3.}
若 \(\mathbf O_0\) 是由
\[
  \mathbf O_0(S)=\mathbb Z
\]
定义的 \(\widehat{\mathscr C}\)-环
(不能将它与常对象 \(\mathbb Z\) 所对应的函子混淆), 则
\(\mathbf O_0\)-模范畴同构于交换 \(\widehat{\mathscr C}\)-群范畴.

\medskip
\noindent\SGAInlineRefTarget{sga3:i:definition:3-1-4}\textbf{定义 3.1.4.}
注意, 若 \(\mathbf F\) 是一个 \(\mathbf O\)-模, 则对每个
\(S\in\operatorname{Ob}(\mathscr C)\), \(\mathbf F_S\) 都是一个
\(\mathbf O_S\)-模. 因此可由
\[
  \underline{\operatorname{Hom}}_{\mathbf O}(\mathbf F,\mathbf F')(S)
    =\operatorname{Hom}_{\mathbf O_S}(\mathbf F_S,\mathbf F'_S).
\]
定义一个阿贝尔 \(\widehat{\mathscr C}\)-群
\(\underline{\operatorname{Hom}}_{\mathbf O}(\mathbf F,\mathbf F')\).
同样可以定义对象
\[
  \underline{\operatorname{Isom}}_{\mathbf O}(\mathbf F,\mathbf F'),
  \qquad
  \underline{\operatorname{End}}_{\mathbf O}(\mathbf F),
  \qquad\text{及}\qquad
  \underline{\operatorname{Aut}}_{\mathbf O}(\mathbf F),
\]
其中最后一个对象在由自同构复合所诱导的结构下是一个
\(\widehat{\mathscr C}\)-群.

\medskip
\noindent\SGAInlineRefTarget{sga3:i:definition:3-2}\textbf{定义 3.2.}
设 \(\mathbf O\) 是一个 \(\widehat{\mathscr C}\)-环,
\(\mathbf F\) 是一个 \(\mathbf O\)-模, 而 \(\mathbf G\) 是一个
\(\widehat{\mathscr C}\)-群.
\(\mathbf F\) 上的 \(\mathbf G\)-\(\mathbf O\)-模结构是指一个
\(\mathbf G\)-对象结构, 使得对每个
\(S\in\operatorname{Ob}(\mathscr C)\) 以及每个
\(g\in\mathbf G(S)\), 由 \(g\) 定义的 \(\mathbf F(S)\) 的自同构
都是其 \(\mathbf O(S)\)-模结构的自同构.

等价地, \SGARefLink{sga3:i:subsection:2-3}{2.3} 中定义的
\(\widehat{\mathscr C}\)-群态射
\[
  \rho:\mathbf G\longrightarrow
    \underline{\operatorname{Aut}}(\mathbf F)
\]
将 \(\mathbf G\) 映入 \(\underline{\operatorname{Aut}}(\mathbf F)\)
的 \(\widehat{\mathscr C}\)-子群
\(\underline{\operatorname{Aut}}_{\mathbf O}(\mathbf F)\).
因此, 在 \(\mathbf O\)-模 \(\mathbf F\) 上赋予一个
\(\mathbf G\)-\(\mathbf O\)-模结构, 等价于给定一个
\(\widehat{\mathscr C}\)-群态射
\[
  \rho:\mathbf G\longrightarrow
    \underline{\operatorname{Aut}}_{\mathbf O}(\mathbf F).
\]
自然可以定义阿贝尔群
\(\operatorname{Hom}_{\mathbf G\text{-}\mathbf O}
  (\mathbf F,\mathbf F')\), 从而得到
\(\mathbf G\)-\(\mathbf O\)-模的加性范畴, 记作
\((\mathbf G\text{-}\mathbf O\text{-Mod.})\).

\medskip
\noindent\SGAInlineRefTarget{sga3:i:remark:3-2-1}\textbf{注 3.2.1.}
读者会注意到, \((\mathbf G\text{-}\mathbf O\text{-Mod.})\)
也可以定义为 \(\mathbf O[\mathbf G]\)-模范畴, 其中
\(\mathbf O[\mathbf G]\) 是 \(\widehat{\mathscr C}\)-群 \(\mathbf G\)
在 \(\widehat{\mathscr C}\)-环 \(\mathbf O\) 上的代数,
其定义是清楚的.
\footnote{编者注: 以下句子为新增内容, 将用于\SGARefLink{sga3:i:section:5}{第 5 节}.}
因此, 由 \SGARefLink{sga3:i:lemma:3-1-1}{3.1.1} 和
\SGARefLink{sga3:i:proposition:3-1-2}{3.1.2},
\((\mathbf G\text{-}\mathbf O\text{-Mod.})\) 是一个满足公理
\(\mathrm{(AB\ 5)}\) 的阿贝尔范畴; 此外, 若 \(\mathscr C\)
是小范畴, 则 \((\mathbf G\text{-}\mathbf O\text{-Mod.})\)
有足够多的内射对象.

本节的所有构造都立即适用于 \(\mathbf G\) (或 \(\mathbf O\),
或两者) 由 \(\mathscr C\) 的一个对象表出的情形,
该对象于是配备了相应的代数结构.

我们已简要讨论了本讨论班后文中出现的主要代数结构.
对于其他结构 (例如 \(\mathbf O\)-Lie 代数结构),
相信本节已为读者提供了足够多的例子, 使读者能在每个具体情形中
运用 \SGARefLink{sga3:i:subsection:2-2}{2.2} 中勾画的一般机制.

现在将以上所做的工作应用于概形范畴, 记作
\((\mathbf{Sch})\), 更一般地, 应用于范畴
\((\mathbf{Sch})_{/S}\), 也记作 \((\mathbf{Sch}_{/S})\).

% End of assigned component. Section 4 begins on Exp1 component/printed p. 21,
% combined-reader p. 31, in the next unit.

% Translation coverage: entire assigned component; no unresolved translation TODOs.
% Statistics: source lines=292; target lines=279; inline math=175; display math=18; align blocks=1; unresolved TODOs=0.
